\chapter{پیاده‌سازی سامانه}\label{Chap:chap4}
{}

\section{مقدمه}

در این فصل، جزئیات پیاده‌سازی مؤلفه‌های مختلف سامانه‌ تشریح می‌گردد. معماری سامانه در فصل قبل معرفی شد و در این فصل، نحوه‌ی پیاده‌سازی هر یک از مؤلفه‌ها شامل آماده‌سازی داده‌ها، پایگاه داده‌ی برداری، عامل مسیریابی، عامل ترجمه، موتور بازیابی، بازرتبه‌بندی، تولید پاسخ با مدل‌زبانی، نقش‌های کاربردی و رابط کاربری به تفکیک بیان می‌شود. تمام مؤلفه‌ها با استفاده از زبان برنامه نویسی پایتون و کتابخانه‌های متن‌باز مربوطه پیاده‌سازی شده‌اند.


به منظور افزایش خوانایی و تمرکز بر مبانی نظری و طراحی، از درج حجم بالای کدهای پیاده‌سازی در متن اصلی پایان‌نامه خودداری شده است. کد‌های کامل و مستند این سامانه، به همراه راهنمای نصب و اجرا، در یک مخزن عمومی بر روی بستر گیت‌هاب قرار داده شده است. ارجاعات به بخش‌های خاص کد، در طول این فصل با ذکر لینک مستقیم به فایل مربوطه در مخزن مذکور ارائه شده‌اند. برای دسترسی به نسخه‌ی دقیق کد‌های مورد استفاده در پژوهش، لطفا به کامیت \texttt{abc123def456} در مخزن زیر مراجعه فرمایید:

\vspace{0.3cm}
\centerline{\texttt{https://github.com/your-username/urology-rag-system}}
\vspace{0.3cm}

\section{پیاده‌سازی آماده‌سازی و پیش‌پردازش داده‌ها}

فرآیند آماده‌سازی داده‌ها شامل استخراج متن، پاکسازی، قطعه‌بندی
\LTRfootnote{chunking}
 و استخراج فراداده
\LTRfootnote{metadata}
 است که در ادامه به تشریح هر یک پرداخته می‌شود.

\subsection{استخراج متن}

متون کتاب‌ها با استفاده از کتابخانه‌ی پی‌ام‌یو‌پی‌دی‌اف۴ال‌ال‌ام
\LTRfootnote{pymupdf4llm}
استخراج می‌شوند. این کتابخانه امکان استخراج متن همراه با ساختار سرفصل‌ها و شناسایی جداول را فراهم می‌کند. کد مربوط به استخراج متن در فایل \texttt{src/data\_extraction.py} در مخزن گیت‌هاب موجود است.

\subsection{پاک‌سازی داده}

پس از استخراج متن، مرحله‌ی پاک‌سازی داده انجام می‌شود. در این مرحله، عملیات زیر بر روی متن اعمال می‌گردد:

\begin{enumerate}
\item حذف کاراکترهای اضافی و فاصله‌های بی‌رویه
\item حذف علائم غیراستاندارد و کاراکترهای کنترلی
\item عادی‌سازی\LTRfootnote{Normalization} یونیکد
\item یکپارچه‌سازی علائم نگارشی
\item مشخص‌سازی سرفصل‌ها 
\item نوشتن متن توصیفی برای تصاویر، نمودار‌ها و برخی جداول
\end{enumerate}

\subsection{قطعه‌بندی سلسله‌مراتبی و حفاظت از جداول}

فرایند قطعه‌بندی متون در دو مرحله‌ی مجزا با استفاده از دو اسکریپت مجزا انجام می‌پذیرد: ابتدا اسکریپت \texttt{parser-into-json.py} و سپس اسکریپت \texttt{chunking.py}. این رویکرد دو مرحله‌ای امکان کنترل دقیق‌تر بر فرایند قطعه‌بندی و حفظ ساختار سلسله‌مراتبی اسناد را فراهم می‌کند.

\subsubsection{مرحله‌ی اول: پردازش ساختاری و تولید قطعات کلان}

در مرحله‌ی اول، اسکریپت \texttt{parser-into-json.py} اجرا می‌شود. وظیفه‌ی اصلی این اسکریپت، تبدیل فایل‌های متنی خام به قطعات کلان ساختاریافته بر اساس سرفصل‌های فصول است. فرایند این مرحله به صورت زیر است:

\begin{enumerate}
	\item \textbf{خوانش فایل‌های متنی}: اسکریپت تمامی فایل‌های متنی با پسوند \texttt{.txt} موجود در پوشه‌ی ورودی را شناسایی و خوانش می‌کند.
	\item \textbf{پاک‌سازی اولیه}: در این گام، عملیات زیر بر روی متن اعمال می‌شود:
	\begin{itemize}
		\item حذف نشانه‌های درشت‌نمایی (\texttt{**}) که از ساختار مارک‌داون باقی مانده‌اند.
		\item یکپارچه‌سازی کلمات شکسته‌شده در انتهای خطوط (De-hyphenation) با استفاده از الگوی عبارت باقاعده.
	\end{itemize}
	\item \textbf{قطعه‌بندی مبتنی بر سرفصل}: با استفاده از کلاس \texttt{MarkdownHeaderTextSplitter} از کتابخانه‌ی \texttt{LangChain}، متن بر اساس الگوی سرفصل‌های سطح فصل (شامل هفت علامت مربع) به قطعات کلان تفکیک می‌شود. پارامتر \texttt{strip\_headers=False} باعث می‌شود تا عنوان هر فصل درون بدنه‌ی متن حفظ شود.
	\item \textbf{تولید فراداده و ذخیره‌سازی}: برای هر قطعه‌ی کلان، یک شناسه‌ی یکتا (\texttt{UUIDv4})، شماره‌ی ترتیب، شناسه‌ی سند، نام فصل، مسیر سلسله‌مراتبی (breadcrumb) و آمارهای مربوط به طول متن (تعداد کاراکتر و تعداد توکن) استخراج می‌شود. خروجی نهایی به صورت یک فایل \texttt{.json} برای هر سند ذخیره می‌گردد.
\end{enumerate}

کد کامل این مرحله در فایل \texttt{parser-into-json.py} در مخزن گیت‌هاب موجود است.

\subsubsection{مرحله‌ی دوم: قطعه‌بندی ریزدانه با حفاظت از جداول}

در مرحله‌ی دوم، اسکریپت \texttt{chunking.py} اجرا می‌شود. این اسکریپت، قطعات کلان تولیدشده در مرحله‌ی قبل را دریافت کرده و با اعمال یک رویکرد دقیق، آن‌ها را به قطعات ریزدانه‌تر تقسیم می‌کند. ویژگی کلیدی این مرحله، \textbf{حفاظت از ساختار جداول} در فرایند قطعه‌بندی است تا از شکسته‌شدن جداول بالینی و از دست رفتن اطلاعات آن‌ها جلوگیری شود. فرایند این مرحله به صورت زیر است:

\begin{enumerate}
	\item \textbf{حفاظت و ماسک‌گذاری جداول (\texttt{Protection Table})}: پیش از قطعه‌بندی، ساختار جداول به همراه کپشن‌های بالادستی آن‌ها (تا ۳ سطر قبل) شناسایی می‌شوند. برای شناسایی جداول، از معیارهایی نظیر تعداد کاراکترهای خط عمودی (|) در هر سطر استفاده می‌شود. همچنین کپشن‌های جداول با استفاده از الگوی عبارت باقاعده \texttt{Table X.X} یا \texttt{eTable} شناسایی می‌گردند. پس از شناسایی، هر جدول با یک کلید موقت یکتا (نظیر \texttt{@@TABLE\_i\_UUID@@}) جایگزین می‌شود تا در فرایند قطعه‌بندی دست‌نخورده باقی بماند.
	
	\item \textbf{قطعه‌بندی بازگشتی با طول متغیر}: متن محافظت‌شده با استفاده از کلاس 
	\texttt{RecursiveCharacterTextSplitter} از کتابخانه‌ی \texttt{LangChain} به قطعات ریزدانه تفکیک می‌شود. پارامترهای تنظیمی این مرحله عبارت‌اند از:
	\begin{itemize}
		\item \textbf{طول بیشینه‌ی قطعه (\texttt{CHUNK\_SIZE})}: برابر با ۱۰۰۰ کاراکتر به منظور پوشش دادن متن یک پاراگراف در یک قطعه.
		\item \textbf{طول همپوشانی (\texttt{CHUNK\_OVERLAP})}: برابر با ۲۰۰ کاراکتر به منظور حفظ پیوستگی مفاهیم بین دو قطعه‌ی متوالی.
	\end{itemize}
	
	\item \textbf{بازنشانی جداول و پاک‌سازی نهایی}: پس از اتمام قطعه‌بندی، کلیدهای موقت با متن کامل جدول جایگزین می‌شوند و متن هر قطعه از فاصله‌های اضافی پاک‌سازی می‌گردد.
	
	\item \textbf{مدیریت قطعات مجزا و ضعیف (\texttt{Chunks Dangling})}: چنانچه قطعه‌ای صرفاً شامل یک تیتر باشد (Heading-only) یا طول آن کمتر از ۱۵۰ کاراکتر باشد، در بافر موقت ذخیره شده و به ابتدای قطعه‌ی معتبر بعدی یا انتهای آخرین قطعه‌ی همان بخش والد الصاق می‌گردد.
	
	\item \textbf{تولید فراداده‌ی نهایی}: برای هر قطعه‌ی ریزدانه، فراداده‌های زیر استخراج و ذخیره می‌شود:
	\begin{itemize}
		\item شناسه‌ی یکتای قطعه (\texttt{UUIDv4})
		\item شناسه‌ی قطعه‌ی والد (برای امکان پیاده‌سازی بازیابی سند والد)
		\item شماره‌ی ترتیب قطعه در کل پیکره
		\item شناسه‌ی سند منبع
		\item مسیر سلسله‌مراتبی (breadcrumb)
		\item متن اصلی قطعه
		\item پرچم وجود جدول در قطعه (has\_table)
	\end{itemize}
\end{enumerate}

کد کامل این مرحله در فایل \texttt{chunking.py} در مخزن گیت‌هاب موجود است.

در نهایت، پس از اعمال مراحل فوق بر روی کل پیکره‌ی متنی، تعداد ۳۶۲۸۵ قطعه (Chunk) ساختاریافته تولید و برای تزریق به پایگاه داده‌ی برداری آماده گردید.

\section{پیاده‌سازی پایگاه داده برداری و نمایه‌سازی}

در سیستم‌های تولید افزوده با بازیابی (\lr{RAG})، پایگاه داده برداری\LTRfootnote{Vector Database} نقش حافظه معنایی بلندمدت را ایفا می‌کند. به منظور تضمین مقیاس‌پذیری بالا، پایداری داده‌ها و دسترسی‌پذیری مداوم، در این پژوهش از پایگاه داده \lr{Qdrant} بر بستر ابری (\lr{Qdrant Cloud}) استفاده شده است. فرایند پیاده‌سازی و نمایه‌سازی\LTRfootnote{Indexing} اسناد در این بخش، با در نظر گرفتن چالش‌های پردازش داده‌های حجیم طراحی شده است که در ادامه مؤلفه‌های کلیدی آن تشریح می‌گردد.

\subsection{استراتژی تعبیه‌سازی ترکیبی}

برای دستیابی به بالاترین دقت در بازیابی اسناد، به‌ویژه در متون تخصصی که واژگان کلیدی اهمیت بالایی دارند، از معماری جستجوی ترکیبی\LTRfootnote{Hybrid Search} استفاده شد. این استراتژی دو مدل بازنمایی متفاوت را به صورت همزمان روی هر قطعه‌متن\LTRfootnote{Chunk} اعمال می‌کند:

\begin{itemize}
	\item \textbf{بردار متراکم\LTRfootnote{Dense Vector}:} برای درک معنایی و مفهومی عمیق متون، از مدل زبانی \lr{\texttt{BAAI/bge-large-en-v1.5}} استفاده شد. این مدل، متون را به بردارهایی با ابعاد ۱۰۲۴ نگاشت می‌کند. برای محاسبه شباهت در فضای برداری، از معیار فاصله کسینوسی\LTRfootnote{Cosine Distance} همراه با نرمال‌سازی بردارها پیش از تزریق به پایگاه داده استفاده شده است.
	\item \textbf{بردار تنک\LTRfootnote{Sparse Vector}:} برای حفظ دقت در تطبیق کلمات کلیدی نادر و تخصصی\LTRfootnote{Lexical Matching}، مدل \lr{\texttt{Splade\_PP\_en\_v1}} به کار گرفته شد. بردارهای تنک تولید شده، روی دیسک\LTRfootnote{On-disk Indexing} نمایه‌سازی می‌شوند تا مصرف حافظه اصلی (\lr{RAM}) در پایگاه داده بهینه گردد.
\end{itemize}

به منظور بهینه‌سازی زمان پردازش، تخصیص منابع سخت‌افزاری در خط لوله\LTRfootnote{Pipeline} تعبیه‌سازی به صورت تفکیک‌شده انجام گرفت؛ به گونه‌ای که پردازش‌های سنگین مدل متراکم به پردازنده گرافیکی (\lr{GPU}) محول شد و استنتاج مدل تنک به صورت موازی بر روی پردازنده مرکزی (\lr{CPU}) صورت پذیرفت.

\subsection{مدیریت فراداده‌ها}

هر نقطه\LTRfootnote{Point} در پایگاه داده، علاوه بر بردارهای متراکم و تنک، شامل اطلاعاتی نظیر شناسه سند اصلی (\lr{\texttt{doc\_id}})، مسیر ساختاری یا عنوان بخش‌ها (\lr{\texttt{breadcrumb}})، نشانگر وجود جدول در متن (\lr{\texttt{has\_table}}) و طول قطعه‌متن می‌باشد. این طراحی، امکان استفاده از فیلترهای پیشرفته فراداده\LTRfootnote{Metadata Filtering} را پیش از انجام جستجوی برداری برای موتور \lr{RAG} فراهم می‌سازد.

\subsection{معماری تزریق داده مقاوم در برابر خطا}

با توجه به حجم بالای قطعه‌متن‌ها و احتمال بروز خطاهای شبکه‌ای در زمان ارتباط با سرویس ابری، اسکریپت تزریق داده\LTRfootnote{Upsertion} با رویکردی مقاوم در برابر خطا طراحی و پیاده‌سازی شد:

\begin{enumerate}

\item \textbf{خاصیت ادامه‌پذیری:} سیستم پیش از آغاز فرایند نمایه‌سازی، تمامی شناسه‌های موجود در فضای ابری را استخراج می‌کند. این مکانیزم تضمین می‌کند که در صورت قطع ارتباط یا توقف فرایند، در اجرای مجدد از تزریق داده‌های تکراری جلوگیری شده و پردازش دقیقاً از نقطه توقف از سر گرفته شود.

\item \textbf{پردازش دسته‌ای و مدیریت حافظه:} فرایند تبدیل و ارسال داده‌ها در قالب دسته‌های ۶۴تایی\LTRfootnote{Batch Processing} انجام می‌پذیرد. همچنین، برای جلوگیری از نشت حافظه\LTRfootnote{Memory Leak} در پردازش‌های طولانی‌مدت بر روی \lr{GPU}، مکانیزم بازیافت حافظه\LTRfootnote{Garbage Collection} پس از اتمام هر دسته به صورت صریح فراخوانی می‌شود.

\item \textbf{ثبت نقاط کنترلی\LTRfootnote{Checkpointing}:} در صورت بروز استثنائات\LTRfootnote{Exceptions} غیرقابل کنترل، سیستم وضعیت فعلی پردازش (شماره اندیس توقف) را در یک فایل کنترلی ثبت می‌کند تا هیچ‌گونه از دست‌رفتگی داده\LTRfootnote{Data Loss} در خط لوله نمایه‌سازی رخ ندهد.

\end{enumerate}

\section{پیاده‌سازی عامل ترجمه در معماری چندزبانه}

با توجه به اینکه بخش عمده‌ای از پایگاه دانش معتبر در حوزه اورولوژی به زبان انگلیسی است، اما کاربران نهایی (پزشکان یا بیماران ایرانی) پرسش‌های خود را به زبان فارسی مطرح می‌کنند، پیاده‌سازی یک معماری RAG چندزبانه\LTRfootnote{Cross-lingual RAG} الزامی بود. برای این منظور، یک «عامل ترجمه» اختصاصی طراحی شد که وظیفه مدیریت جریان تبدیل زبان در دو انتهای مسیر پردازش را بر عهده دارد.

\subsection{طراحی زنجیره‌های ترجمه و مهندسی پرامپت}
برای مدیریت ترجمه‌ها از فریم‌ورک LangChain و ابزار \lr{ChatPromptTemplate} استفاده شد. معماری این عامل به دو زنجیره\LTRfootnote{Chain} مجزا تقسیم می‌شود:

\textbf{۱. ترجمه فارسی به انگلیسی (همراه با بهینه‌سازی پرس‌وجو):}
چالش اصلی در ترجمه پرسش کاربر، از دست رفتن کلیدواژه‌های تخصصی و کاهش دقت در مرحله جستجوی تشابه\LTRfootnote{Similarity Search} در پایگاه داده برداری است. برای حل این مشکل، پرامپت\LTRfootnote{Prompt} سیستم به گونه‌ای مهندسی شد که مدل زبانی علاوه بر ترجمه دقیق متن، در صورت تشخیص ماهیت بالینی پرسش، کلیدواژه‌های استاندارد پزشکی، نام روش‌های درمانی و مراحل بیماری (Staging) مرتبط را به انتهای ترجمه اضافه کند. این تکنیک که به عنوان «بسط پرس‌وجو»\LTRfootnote{Query Expansion/Optimization} شناخته می‌شود، غنای معنایی بردار تولید شده را به شدت افزایش می‌دهد.
همچنین قوانین سخت‌گیرانه‌ای برای جلوگیری از اضافه‌گویی مدل (مانند پرهیز از ارائه پاسخ مستقیم به جای ترجمه) و الزام به نویسه‌گردانی\LTRfootnote{Transliteration} دقیق نام داروها (مانند تبدیل تامسولوسین به \texttt{Tamsulosin}) در دستورالعمل سیستم تعبیه شد.

\textbf{۲. ترجمه انگلیسی به فارسی (تولید خروجی نهایی):}
پس از بازیابی اطلاعات و تولید پاسخ نهایی توسط مدل پایه به زبان انگلیسی، این زنجیره وظیفه تبدیل پاسخ به زبان فارسی رسمی و روان را بر عهده دارد. در این مرحله، مدل اکیداً از ابداع واژگان غیرمعمول (Hallucination) منع شده و موظف است خروجی را بدون مقدمه‌چینی اضافی ارائه دهد.

\subsection{یکپارچه‌سازی با گراف جریان‌کار}
عامل ترجمه به عنوان دو گره\LTRfootnote{Node} مستقل در معماری LangGraph پیاده‌سازی شد.

گره 
\lr{\texttt{translate\_to\_english\_node}} در ابتدای گراف، پرسش فارسی کاربر را از وضعیت سیستم (\lr{\texttt{AgenticRAGState}}) دریافت کرده و نسخه انگلیسی و بهینه‌شده آن را به عنوان \lr{\texttt{english\_query}} در وضعیت به‌روزرسانی می‌کند. در انتهای گراف نیز، گره \lr{\texttt{translate\_to\_persian\_node}} پاسخ نهایی استخراج‌شده (\lr{\texttt{rag\_response}}) را دریافت و به فارسی برمی‌گرداند.

\subsection{ارزیابی و اعتبارسنجی ماژول ترجمه}
برای اطمینان از صحت عملکرد عامل ترجمه، یک اسکریپت آزمون جامع توسعه یافت که سناریوهای زیر را مورد ارزیابی قرار می‌دهد:
\begin{itemize}
    \item \textbf{صحت ترجمه نام داروها:} بررسی توانایی سیستم در تبدیل دقیق داروهای رایج اورولوژی (مانند فیناستراید، دوتاستراید و سیپروفلوکساسین) به معادل استاندارد انگلیسی جهت جلوگیری از خطای بازیابی.
    \item \textbf{حفظ اصطلاحات بالینی:} ارزیابی ترجمه علائم بیماری (مانند تکرر ادرار) و مفاهیم تشخیصی (مانند هیپرپلازی خوش‌خیم پروستات).
    \item \textbf{استحکام ساختاری:} اطمینان از عدم تولید متون زائد و اضافه‌گویی توسط مدل.
\end{itemize}
نتایج این آزمون‌ها نشان داد که اعمال محدودیت‌های سخت‌گیرانه در پرامپت‌ها و استفاده از پارسرهای خروجی (\lr{\texttt{StrOutputParser}})، پایداری سیستم را در پردازش زبان طبیعی بهبود بخشیده و آن را برای محیط‌های عملیاتی حساس آماده ساخته است.


\section{طراحی و پیاده‌سازی موتور جستجوی ترکیبی (\lr{\texttt{Hybrid Retriever}})}

مرحله بازیابی اطلاعات در سامانه‌های RAG پزشکی با چالش‌های دوگانه‌ای مواجه است. از یک سو، سیستم باید قادر به درک معنایی پرسش‌های طولانی و پیچیده بیماران باشد، و از سوی دیگر، باید بتواند کلمات کلیدی دقیق مانند نام داروها، ژن‌ها، یا گونه‌های باکتریایی (نظیر \lr{Proteus} و \lr{Klebsiella}) را با دقت بالا تطبیق دهد. برای حل این تعارض، یک معماری «جستجوی ترکیبی»\LTRfootnote{Hybrid Search} بر بستر پایگاه داده برداری Qdrant طراحی و پیاده‌سازی شد.

\subsection{مدل‌های تعبیه‌سازی دوگانه}
کلاس \lr{\texttt{MedicalHybridRetriever}} بر اساس تلفیق دو رویکرد متفاوت در تولید بردارهای ویژگی عمل می‌کند:

\begin{itemize}
    \item \textbf{تعبیه‌سازی متراکم (Dense Embedding):} برای درک بافتار و شباهت معنایی، از مدل پیش‌آموزش‌دیده \lr{BAAI/bge-large-en-v1.5} استفاده شد. این مدل به دلیل عملکرد بالا در بنچمارک‌های MTEB، قادر است ارتباطات پنهان در متون پزشکی را استخراج کند. پردازش این بخش برای بهینه‌سازی سرعت، به صورت پویا روی شتاب‌دهنده سخت‌افزاری (GPU) نگاشت می‌شود.
    \item \textbf{تعبیه‌سازی تُنُک (Sparse Embedding):} برای جبران ضعف مدل‌های متراکم در تطبیق دقیق کلمات کلیدی (Exact Match)، مدل \lr{prithivida/Splade\_PP\_en\_v1} از کتابخانه \lr{fastembed} به کار گرفته شد. بردارهای تُنُک مشابه الگوریتم BM25 عمل کرده اما وزن‌دهی کلمات را بر اساس یادگیری عمیق بهینه‌سازی می‌کنند.
\end{itemize}

\subsection{استراتژی پیش‌واکشی و تلفیق رتبه‌ها (RRF)}
یکی از نوآوری‌های مهندسی در این بخش، مدیریت نحوه کوئری زدن به پایگاه داده است. از آنجا که مقیاس امتیازات (Scores) در فضای برداری متراکم (عموماً شباهت کسینوسی) با فضای برداری تُنُک متفاوت است، امکان جمع مستقیم آن‌ها وجود ندارد. 

برای حل این مشکل، از مکانیزم «تلفیق رتبه‌های متقابل»\footnote{Reciprocal Rank Fusion (RRF)} استفاده شد که مستقیماً در موتور Qdrant پشتیبانی می‌شود. فرآیند بازیابی طی یک استراتژی بهینه‌سازی‌شده به شکل زیر انجام می‌گردد:
\begin{enumerate}
    \item \textbf{پیش‌واکشی نامتقارن\footnote{Asymmetric Prefetching}:} سیستم ابتدا تعداد $2 \times K$ سند را از طریق جستجوی متراکم و تعداد $K$ سند را از طریق جستجوی تُنُک پیش‌واکشی (Prefetch) می‌کند. واکشی دو برابری در فضای متراکم، تضمین می‌کند که اسناد دارای ارتباط معنایی بالا، فرصت حضور در مرحله رتبه‌بندی نهایی را از دست ندهند.
    \item \textbf{اعمال RRF:} نتایج دو استخر واکشی‌شده بر اساس رتبه آن‌ها (و نه امتیاز خام برداری) با یکدیگر ادغام و رتبه‌بندی مجدد می‌شوند.
\end{enumerate}

\subsection{حفظ بافتار سلسله‌مراتبی در نتایج}
همان‌طور که در بخش پیش‌پردازش داده‌ها اشاره شد، قطعات متنی به تنهایی ممکن است معنای خود را از دست بدهند. لذا در مرحله استخراج نتایج نهایی (\lr{Payload Extraction}) از پایگاه داده، علاوه بر متن اصلی و شناسه قطعه (\lr{\texttt{chunk\_id}})، فراداده \lr{\texttt{breadcrumb}} نیز استخراج شده و به همراه نتیجه بازگردانده می‌شود. این امر به مدل زبانی تولیدکننده پاسخ (LLM) کمک می‌کند تا جایگاه دقیق اطلاعات را در فصل و بخش مربوطه از کتاب مرجع پزشکی درک کند.

\subsection{اعتبارسنجی با شرح‌حال بالینی}
به منظور ارزیابی عملکرد موتور جستجوی ترکیبی، سیستم با استفاده از شرح‌حال‌های بالینی\LTRfootnote{Clinical Vignettes} پیچیده مورد تست قرار گرفت. به عنوان نمونه، پرسشی شامل علائم یک بیمار زن ۵۸ ساله با سابقه عفونت‌های مکرر، نتایج سونوگرافی کلیه (هیدرونفروز) و کشت ادرار وارد سیستم شد. نتایج نشان داد که بازیاب ترکیبی با موفقیت توانست به دلیل وجود کلمات کلیدی تخصصی و درک معنایی از وضعیت بیمار، مرتبط‌ترین متون را جهت پاسخ به سؤالات تشخیص افتراقی \lr{\texttt{Differential Diagnosis}} استخراج نماید.

\subsection{پیاده‌سازی بازرتبه‌بندی}
\subsection{پیاده‌سازی تولید پاسخ با مدل زبانی}
\section{پیاده‌سازی نقش‌های پزشک و بیمار}
\section{معماری لایه سرویس و رابط کاربری}

به منظور برقراری ارتباط میان هسته‌ی پردازشی سیستم (موتور RAG) و کاربران نهایی، لایه‌ی واسط برنامه‌نویسی\footnote{API} سامانه با استفاده از چارچوب فست‌ای‌پی‌آی\footnote{FastAPI} طراحی و پیاده‌سازی شده است. معماری این بخش مبتنی بر الگوی طراحی ماژولار بوده و با هدف جداسازی دغدغه‌ها (Separation of Concerns)، وظایف مختلف اعم از اعتبارسنجی داده‌ها، مسیریابی درخواست‌ها و تولید پاسخ از یکدیگر تفکیک شده‌اند.

\subsection{طراحی و پیاده‌سازی لایه API}

به جای توسعه‌ی یک رابط برنامه‌نویسی با ساختار سفارشی، استراتژی معماری در این سامانه بر پایه‌ی تطابق کامل با استاندارد API شرکت OpenAI بنا شده است. این تصمیم مهندسی موجب می‌شود تا سامانه بتواند به صورت «وصل و اجرا» (Plug-and-Play) با تمامی رابط‌های کاربری و ابزارهای توسعه‌یافته برای اکوسیستم OpenAI یکپارچه شود. مؤلفه‌های اصلی این لایه عبارت‌اند از:

\subsubsection{مدیریت طرح‌واره‌ها و اعتبارسنجی (Data Schemas)}

برای تضمین یکپارچگی داده‌های ورودی و خروجی، از مدل‌سازی داده مبتنی بر کتابخانه‌ی \texttt{Pydantic} استفاده شده است. تمامی درخواست‌های چت (مانند \texttt{ChatCompletionRequest}) و ساختار پیام‌ها، پیش از ورود به لایه‌ی پردازش، به صورت خودکار اعتبارسنجی می‌شوند. این لایه وظیفه دارد اطمینان حاصل کند که فرمت درخواست‌ها دقیقاً با ساختار استاندارد OpenAI هم‌خوانی دارد؛ در غیر این صورت، خطاهای معنادار در سطح API بازگردانده می‌شود.

\subsubsection{مسیریابی و مدیریت نقاط پایانی (Endpoints)}

هسته‌ی اصلی ارتباطی سامانه در دو نقطه‌ی پایانی (Endpoint) کلیدی متمرکز شده است:

\begin{itemize}

\item \textbf{نقطه‌ی پایانی \texttt{/v1/models}:} این مسیر وظیفه‌ی معرفی مدل‌های در دسترس را بر عهده دارد. در یک رویکرد ابتکاری، به جای معرفی مدل‌های زبانی پایه (مانند LLaMA یا GPT)، «پرسونای پاسخ‌گو» به عنوان مدل به رابط کاربری معرفی می‌شود. بر این اساس، دو شناسه مدل به نام‌های \texttt{urology-rag-doctor} (با لحن تخصصی پزشکی) و \texttt{urology-rag-patient} (با لحن ساده و بیمارمحور) در این نقطه تعریف شده‌اند تا سامانه‌های بالادستی بتوانند آن‌ها را شناسایی کنند.

\item \textbf{نقطه‌ی پایانی \texttt{/v1/chat/completions}:} این مسیر دریافت‌کننده‌ی اصلی پرسش‌های کاربران است. منطق پیاده‌سازی‌شده در این بخش، پس از دریافت درخواست، ابتدا آخرین پیام کاربر را استخراج کرده و سپس بر اساس مدل انتخابی (شناسه پزشک یا بیمار)، نقش کاربر را برای تنظیم لحن سیستم (System Prompt) مشخص می‌کند. در نهایت، درخواست به لایه‌ی ارکستراسیون (Orchestration Layer) ارسال می‌گردد.

\end{itemize}

\subsubsection{ارکستراسیون جریان پردازش (LangGraph Service)}

مدیریت جریان پیچیده‌ی RAG — شامل بازیابی اطلاعات، بازرتبه‌بندی (Reranking) و تولید نهایی پاسخ — بر عهده‌ی یک گراف حالت‌گذار است که با استفاده از کتابخانه‌ی \texttt{LangGraph} پیاده‌سازی شده است. سرویس تولید پاسخ، پرسش کاربر و نوع پرسونای درخواستی را به عنوان ورودی به این گراف تزریق می‌کند. استفاده از LangGraph این امکان را فراهم کرده است تا مراحل پردازش به صورت گره‌های (Nodes) مستقل تعریف شده و در صورت نیاز به افزودن منطق‌های شرطی (مانند بررسی کیفیت بازیابی پیش از تولید پاسخ)، جریان کار به راحتی توسعه‌پذیر باشد.

(جزئیات کامل کدها و ساختار ماژول‌ها در مخزن گیت‌هاب پروژه قابل بررسی است.)

\subsection{یکپارچه‌سازی با رابط کاربری (OpenWebUI)}

برای تعامل بصری کاربران با سامانه، از پلتفرم متن‌باز اوپن‌دبلیو‌یو‌آی\footnote{OpenWebUI} به عنوان رابط کاربری نهایی استفاده شده است.

\subsubsection{مکانیسم اتصال و شناسایی}

از آنجا که لایه‌ی API سامانه به طور کامل با استاندارد OpenAI سازگار شده است، اتصال OpenWebUI به موتور RAG بدون نیاز به هیچ‌گونه تغییر در کدهای سمت کلاینت (Client-side) انجام می‌پذیرد. تنها با تنظیم آدرس پایگاه API (Base URL) به آدرس سرور میزبان این سامانه، OpenWebUI درخواست‌های شناسایی مدل را به مسیر \texttt{/v1/models} ارسال کرده و دو پرسونای تخصصی و عمومی اورولوژی را به عنوان گزینه‌های قابل انتخاب در منوی خود به کاربر نمایش می‌دهد.

\subsubsection{استراتژی ارسال پاسخ}

در معماری فعلی، استراتژی بازگشت پاسخ بر مبنای حالت غیرجریانی\footnote{Non-streaming} تنظیم شده است؛ بدین معنا که خروجی نهایی پس از تکمیل فرایند پردازش در موتور RAG به صورت یکپارچه به رابط کاربری ارسال می‌شود. با این وجود، به دلیل استفاده از قابلیت‌های ناهمگام (Asynchronous) در FastAPI و ساختار استاندارد پیام‌ها، زیرساخت لازم برای افزودن قابلیت پاسخ‌دهی توکن‌به‌توکن (Streaming Response) در توسعه‌های آتی سیستم کاملاً فراهم است.